% ============================================================
% INTEGRAZIONI CAPITOLO 1
% Integrazione framework manageriali — osservazioni Prof. Sasso
% ============================================================
%
% Questo file contiene SOLO i blocchi da inserire nel Capitolo 1.
% Non modificare il resto del capitolo.
%
% ============================================================


% ============================================================
% INTEGRAZIONE 1 — §1.1
% INSERIRE DOPO \subsection{Verso una PA proattiva: il ruolo emergente dell'IA}
% (ossia dopo il paragrafo che termina con "...ragioni fondamentali per cui
%  un modello di deployment dell'IA in ambiente on-premise può offrire
%  vantaggi strategici alla PA.")
% e PRIMA di \section{Piano Triennale ICT 2024--2026}
% ============================================================

\subsection{Chiavi di lettura manageriali per la trasformazione digitale}
\label{subsec:chiavi-manageriali}

L'analisi condotta nelle pagine precedenti ha ricostruito il percorso evolutivo che ha condotto la PA europea dal paradigma dell'\textit{e-Government} a quello, più maturo, del \textit{Digital Government}. Tale percorso, tuttavia, rischia di essere letto in modo prevalentemente normativo e descrittivo se non viene inquadrato attraverso le lenti concettuali offerte dalla letteratura manageriale e degli \textit{Information Systems}. Le scelte di trasformazione digitale --- inclusa quella, centrale per questa tesi, relativa al modello di \textit{deployment} dell'intelligenza artificiale --- non sono mere decisioni tecnologiche: sono decisioni organizzative, strategiche e istituzionali che possono essere analizzate con strumenti teorici consolidati. Il presente paragrafo introduce tre \textbf{framework interpretativi complementari} che costituiranno il filo conduttore analitico dei capitoli successivi.

Il primo è il \textbf{\textit{Technology-Organization-Environment} framework} (\textit{TOE}), sviluppato da \textcite{TornatzkyFleischer1990} per spiegare i processi di adozione tecnologica nelle organizzazioni. Il modello postula che l'adozione di un'innovazione — inclusa quella di sistemi di IA — sia simultaneamente condizionata da tre contesti interrelati: il contesto \textit{tecnologico} (maturità, complessità, compatibilità con i sistemi esistenti e disponibilità dell'infrastruttura), il contesto \textit{organizzativo} (dimensione dell'ente, disponibilità di risorse, competenze del personale, cultura organizzativa e disponibilità della leadership) e il contesto \textit{ambientale} (pressione regolamentare, pressione istituzionale e disponibilità di supporto esterno). La forza del \textit{TOE framework} risiede nella sua capacità di spiegare perché organizzazioni apparentemente simili per dimensione e settore possano presentare traiettorie di adozione radicalmente diverse: la variabile discriminante non è mai il solo fattore tecnologico, ma la combinazione dei tre contesti. Applicato alla PA, il \textit{TOE} permette di analizzare la scelta del modello \textit{on-premise} non come una preferenza tecnica, ma come l'esito di una configurazione specifica di vincoli organizzativi (competenze interne, budget CAPEX disponibile), tecnologici (infrastruttura esistente, compatibilità con i sistemi informativi dell'ente) e ambientali (obblighi normativi come il GDPR e la Strategia Cloud Italia, pressione istituzionale da parte di AgID e ACN).

Il secondo framework è la \textbf{\textit{Institutional Theory}} e, in particolare, il concetto di \textbf{isomorfismo istituzionale} elaborato da \textcite{DiMaggioPow1983}. Gli autori identificano tre meccanismi attraverso cui le organizzazioni tendono ad omologarsi alle norme e alle aspettative del proprio campo istituzionale: l'isomorfismo \textit{coercitivo}, generato da pressioni normative e regolatorie (nel contesto della PA: l'AI Act, la NIS~2, il Perimetro di Sicurezza Nazionale Cibernetica); l'isomorfismo \textit{mimetico}, prodotto dall'incertezza e dall'imitazione di organizzazioni ritenute di successo (la tendenza delle PA italiane a replicare le scelte infrastrutturali di enti «pionieri» o le \textit{best practice} europee del PSN); e l'isomorfismo \textit{normativo}, derivante dai processi di professionalizzazione e dalla diffusione di standard condivisi attraverso reti di esperti, comunità di pratica e programmi di formazione (il ruolo di AgID Academy, della CRUI e dei network di competenza digitale). L'\textit{Institutional Theory} offre uno strumento prezioso per interpretare non soltanto le ragioni dell'adozione dell'IA nella PA, ma anche le resistenze e le asimmetrie che la caratterizzano.

Il terzo framework è la teoria delle \textbf{\textit{Dynamic Capabilities}}, proposta da \textcite{Teece2007} nel solco della \textit{Resource-Based View}. L'autore definisce le \textit{dynamic capabilities} come la capacità di un'organizzazione di integrare, costruire e riconfigurare le proprie competenze interne ed esterne per rispondere — e anticipare — i cambiamenti ambientali. Il framework si articola in tre micro-fondamenti: \textbf{\textit{sensing}} (la capacità di identificare e valutare opportunità e minacce nel contesto esterno), \textbf{\textit{seizing}} (la capacità di mobilitare risorse per cogliere le opportunità identificate) e \textbf{\textit{transforming}} (la capacità di riconfigurare continuamente l'organizzazione per sostenere il vantaggio competitivo nel tempo). Applicato alla PA, il framework delle \textit{dynamic capabilities} permette di distinguere le amministrazioni capaci di governare attivamente la trasformazione digitale da quelle che vi si adattano passivamente o la subiscono: le prime sviluppano capacità di \textit{sensing} dell'ecosistema tecnologico e regolatorio, capacità di \textit{seizing} nella selezione e acquisizione delle soluzioni più adeguate, e capacità di \textit{transforming} nella riorganizzazione dei processi interni attorno alle nuove tecnologie.

Questi tre framework non sono alternativi, ma \textbf{complementari}: il \textit{TOE} illumina i \textit{fattori contestuali} che condizionano l'adozione, la \textit{Institutional Theory} spiega le \textit{dinamiche di campo} che ne orientano le direzioni, le \textit{Dynamic Capabilities} analizzano le \textit{capacità organizzative} che ne determinano il successo. Nel Capitolo~2, ciascuno di questi framework sarà applicato in modo sistematico all'analisi dell'adozione dell'IA nella PA italiana, con particolare attenzione alle implicazioni per la scelta del modello di \textit{deployment on-premise}.


% ============================================================
% INTEGRAZIONE 2 — §1.5
% INSERIRE nel corpo del paragrafo \subsection{Implicazioni per la ricerca}
% (§1.5, ultimo paragrafo del capitolo)
% AGGIUNGERE le seguenti 3 frasi DOPO il punto:
%   "...senza dipendere da API esterne e da connettività verso data center remoti."
% e PRIMA del punto finale che inizia con:
%   "Il capitolo successivo approfondirà..."
% ============================================================

% Le tre frasi da inserire (integrarle nel testo corrente, dopo la frase
% "...senza dipendere da API esterne e da connettività verso data center remoti."):

Sotto questo profilo, gli indicatori DESI e il \textit{Digital Government Index} possono essere riletti non soltanto come strumenti di monitoraggio comparativo, ma come meccanismi istituzionali di \textbf{\textit{sensing}} nel senso di \textcite{Teece2007}: essi segnalano alle PA le opportunità e i ritardi del proprio ecosistema digitale, orientando le decisioni di investimento e le priorità strategiche. Una PA dotata di \textit{dynamic capabilities} sviluppate è in grado di trasformare i dati DESI e DGI in segnali strategici --- identificando, ad esempio, il divario nella dimensione \textit{proactiveness} come un'opportunità specifica per investire in sistemi di IA predittiva --- piuttosto che trattarli come indicatori di \textit{compliance} amministrativa. In questo senso, la misurazione della maturità digitale non è un esercizio puramente valutativo, ma un atto organizzativo con implicazioni dirette sulla capacità delle amministrazioni di \textit{cogliere} (\textit{seizing}) le opportunità offerte dall'intelligenza artificiale.
